\section{Results}
\label{sec:results}

The experiments have been specified, but their artifacts are not yet frozen.
The following placeholders reserve the evidence required to answer each
research question. A placeholder is replaced automatically when its canonical
artifact is placed in the paper's \texttt{figures/} or \texttt{tables/}
directory.

\subsection{RQ1: Representation Coverage and Composition}

\paragraph{Answer placeholder.}
First report which kernel feature maps have faithful expressions, which have
only bounded proxies, and which require a new primitive. Independent parity is
mandatory for every faithful claim. Then illustrate how atomic expressions are
combined into additive, ordered, and relational feature programs. Their induced
discrimination profiles are consequences of the declared features, not
discoveries about what the programs were designed to detect.

\begin{figure}[t]
  \centering
  \artifactfigure{fig01-structural-views.pdf}{Matched graphs, contrasting
  structural views, and mapped witnesses of the intervened structure.}
  \caption{One matched graph pair under several structural views. The final
  figure must show a collapse, a separation, and the overlap between the mapped
  witness and the known intervention.}
  \label{fig:structural-views}
\end{figure}

\begin{table}[t]
  \centering
  \caption{Graph-kernel feature families and their status in the Abstract Graph
  language. ``Faithful'' requires independent parity; a bounded proxy is named
  by its restriction rather than presented as the original kernel.}
  \label{tab:baseline-expressions}
  \IfFileExists{../tables/tab01-baseline-expressions.tex}{%
    \input{../tables/tab01-baseline-expressions.tex}%
  }{%
    \small
    \begin{tabular}{p{0.15\linewidth}p{0.25\linewidth}p{0.23\linewidth}p{0.16\linewidth}}
      \toprule
      Family & Candidate AG expression & Required identity & Current status \\
      \midrule
      Vertex/edge & \texttt{node}/\texttt{edge} & attributed object & parity pending \\
      Shortest path & node pairs by distance & endpoint labels + distance & parity pending \\
      Graphlet & \texttt{graphlet}$(r,k)$ & attributed topology & bounded proxy \\
      WL subtree & neighborhood hashing & rooted local identity & WL-like proxy; exact pending \\
      NSPDK & neighborhoods + distance pairs & rooted neighborhood pair & parity pending \\
      \bottomrule
    \end{tabular}%
  }
\end{table}

\begin{figure}[t]
  \centering
  \artifactfigure{fig02-discrimination-atlas.pdf}{Factor-specific intrinsic
  discrimination for atomic, baseline, concatenated, and composed expressions.}
  \caption{Structural discrimination atlas over predeclared interventions and
  held-out generator settings. Rows are expressions and columns are structural
  factors; no aggregate score replaces the factor-specific profile.}
  \label{fig:discrimination-atlas}
\end{figure}

\begin{figure}[t]
  \centering
  \artifactfigure{fig03-composition-transitions.pdf}{Pairwise transitions for
  graph-valued composition relative to atomic, lossless-concatenation, and flat
  relational controls, with mapped witnesses for representative changes.}
  \caption{Changes in graph-pair discrimination under composition. The
  comparison separates distinctions created by ordered or relational programs
  from those retained by flat feature concatenation, and checks whether an
  equally informed flat materialization recovers the same distinctions.}
  \label{fig:composition-transitions}
\end{figure}

\subsection{RQ2: Discrimination--Complexity Trade-off}

\paragraph{Answer placeholder.}
Report factor-specific non-dominated configurations for runtime, peak memory,
and representation size, with atomic, concatenated, and equally informed flat
controls in every panel. State whether graph-valued composition changes cost or
localization at matched discrimination. Retain dominated and failed
configurations where they show that added structural machinery produces little
or no useful discrimination.

\begin{figure}[t]
  \centering
  \artifactfigure{fig04-pareto-frontiers.pdf}{Factor-specific discrimination
  against runtime, peak memory, and representation size for matched atomic,
  concatenated, flat-materialized, and graph-composed configurations.}
  \caption{Discrimination--complexity Pareto frontiers. Cost dimensions remain
  separate, and aggregate discrimination is secondary to the structural-factor
  panels.}
  \label{fig:pareto-frontiers}
\end{figure}

\begin{table}[t]
  \centering
  \caption{Representative expression trade-offs. Atomic, concatenated, and
  flat-materialized, and graph-composed configurations are compared at matched
  settings.}
  \label{tab:expression-tradeoffs}
  \IfFileExists{../tables/tab02-expression-tradeoffs.tex}{%
    \input{../tables/tab02-expression-tradeoffs.tex}%
  }{%
    \small
    \begin{tabular}{lrrrrl}
      \toprule
      Configuration & Discrimination & Time & Memory & Size & Pareto \\
      \midrule
      Atomic & --- & --- & --- & --- & pending \\
      Concatenated & --- & --- & --- & --- & pending \\
      Flat relational & --- & --- & --- & --- & pending \\
      Graph-composed & --- & --- & --- & --- & pending \\
      \bottomrule
    \end{tabular}%
  }
\end{table}
